\section{A segunda camada de evidência: hierarquia e eficiência}
\label{sec:fase2}

A segunda fase da estimação atacou a mesma pergunta com três famílias
de métodos ainda não usadas no documento: modelos
\textbf{hierárquicos} (que tratam os 11 anos de cada hospital como
parentes, não como observações soltas), a análise envoltória de dados
(\textbf{DEA}, que constrói uma fronteira de melhores práticas e mede
a distância de cada hospital até ela) e a fronteira estocástica
(\textbf{SFA}, a contraparte paramétrica do DEA). O pacote completo de
resultados, com scripts, diagnósticos e um relatório de execução
auditável, está na pasta resultados\_fase2 do repositório; aqui ficam
os achados que mudam ou reforçam a leitura do documento.

\textbf{O que a camada hierárquica confirma.} Reestimados com um
efeito próprio para cada hospital, o tempo de permanência e o custo
real reproduzem quase exatamente os resultados da
Seção~\ref{sec:resultados} (queda de 10,2\% no TMP; custo sem
diferença), a produção e a ocupação continuam maiores nas OSS com
magnitudes menores (33\% e 13\%, porque parte do contraste entre as
redes é absorvida pelo efeito de hospital), e a mortalidade aparece
0,255 menor no logito, com precisão maior que a da
Seção~\ref{sec:resultados}. O teste decisivo, porém, é outro:
\textbf{retirando os cinco conversores, a diferença de mortalidade
entre as redes desaparece} (0,09 positivo, nulo), enquanto TMP,
produção, ocupação e eficiência sobrevivem intactos. É a quantificação
do aviso da Seção~\ref{sec:inferenciafinal}: a mortalidade menor é uma
história dos hospitais que mudaram de gestão, não uma diferença entre
as redes OSS e Direta.

\begin{figure}[H]
\centering
\includegraphics[width=.9\textwidth]{figB02_encolhimento_categorias.png}
\caption{Efeitos de categoria na mortalidade com e sem partial
pooling: com apenas 3 hospitais, o efeito aparente do Privado encolhe
em quase 90\% quando o modelo pondera a evidência pelo tamanho do
grupo.}
\label{fig:encolhimento}
\end{figure}

\comoler{cada linha compara duas estimativas do efeito da mesma
categoria: o círculo escuro trata cada categoria isoladamente; o
quadrado verde deixa as categorias emprestarem informação umas às
outras, na proporção da quantidade de dados de cada uma. Quanto menor
o grupo, maior o recuo do quadrado em direção ao zero: o Privado, com
3 hospitais, quase não sustenta efeito próprio. É a ilustração formal
da regra deste documento de nunca interpretar essa categoria.}

\textbf{Eficiência: a ressalva vem antes do resultado.} O DEA e a SFA
medem eficiência produtiva: quanto produto (saídas, incluída a alta
complexidade) cada hospital extrai dos seus insumos. Um dos insumos é
o \textbf{valor total real} pago pelo SUS, e ele exige a ressalva de
primeira ordem desta seção: o valor do SIH é reembolso por produção
faturada, não medida direta de recursos consumidos, e os regimes de
pagamento diferem por categoria (as OSS recebem pelo contrato de
gestão, de modo que o faturamento por AIH pode divergir do recurso
efetivo; os filantrópicos convivem com subfinanciamento e subsídio
cruzado). Parte da distância entre categorias pode refletir regime de
pagamento, não gestão. Duas verificações sustentam que a ordenação
não é artefato do insumo financeiro: a \textbf{produção física por
leito}, que não usa dinheiro nenhum, já ordena as redes do mesmo jeito
(54 saídas por leito nas OSS e 51 nos Municipais, contra 35 na Direta
e 34 nos Filantrópicos); e o DEA reestimado \textbf{só com leitos como
insumo} preserva e até amplia a vantagem (medianas de 0,68 nas OSS
contra 0,40 na Direta).

Com a ressalva posta, o resultado: na eficiência corrigida de viés
(bootstrap de Simar e Wilson, 2.000 réplicas por ano), as medianas do
período são \textbf{OSS 0,72, Público Municipal 0,64, Direta 0,59 e
Filantrópico 0,58}. Na regressão da distância à fronteira sobre
categoria, complexidade, porte e ano, a vantagem OSS equivale a cerca
de \textbf{40\% menos distância à fronteira que a Direta}, estável
retirando o biênio pandêmico, os cinco conversores ou restringindo ao
suporte comum. Ao contrário da mortalidade, \textbf{a eficiência
maior das OSS é um traço da rede}, presente na comparação pura entre
hospitais distintos. E o receio de os 3 hospitais Privados
distorcerem a fronteira não se confirma: retirados, os escores dos
demais praticamente não se movem (correlação de ordem 1,000), porque
o perfil deles opera longe da fronteira.

\begin{figure}[H]
\centering
\subfig{figE01_violino_te_categoria.png}{Eficiência DEA corrigida de
viés}
\subfig{figS01_te_categoria.png}{Eficiência técnica da fronteira
estocástica}
\caption{Eficiência por categoria nos dois métodos. No painel (b), a
concentração das OSS junto ao valor 1 é uma degeneração numérica do
modelo (o texto explica) e o valor delas não deve ser citado; a ordem
das demais categorias é a informação válida.}
\label{fig:eficiencia}
\end{figure}

\comoler{como nos violinos das seções anteriores: a silhueta engorda
onde há mais hospitais e o número é a mediana do grupo. Compare a
posição dos violinos de OSS e Municipal com os de Direta e
Filantrópico nos dois painéis: métodos de naturezas opostas ordenam as
redes do mesmo jeito. O Privado aparece só como registro.}

\textbf{A contraparte paramétrica e a convergência entre métodos.} A
fronteira estocástica (Translog, preferida ao Cobb Douglas por larga
margem; 94\% da variância composta atribuída à ineficiência) confirma
a ordem das redes, com uma honestidade a registrar: as OSS ficam tão
coladas à fronteira paramétrica que o coeficiente delas na equação de
ineficiência perde identificação (o análogo da separação perfeita da
Seção~\ref{sec:estrategia}), e por isso o NÍVEL da eficiência das OSS
na SFA não é citável; vale a ordem, não o número. O resultado síntese
da fase é a concordância entre os dois métodos: correlação de ordem
de \textbf{0,63 entre as eficiências do DEA e da SFA} (0,65 excluindo
as OSS, mostrando que a concordância não depende do grupo problemático),
crescendo ao longo do período.

\begin{figure}[H]
\centering
\includegraphics[width=.9\textwidth]{figS02_dispersao_sfa_dea.png}
\caption{Convergência entre os dois métodos de eficiência: cada ponto
é um hospital num ano, posicionado pelo escore DEA (horizontal) e pela
eficiência SFA (vertical).}
\label{fig:convergencia}
\end{figure}

\comoler{quanto mais a nuvem sobe da esquerda para a direita, mais os
dois métodos concordam sobre quem é eficiente. A faixa verde colada no
topo é a degeneração das OSS descrita no texto (o modelo paramétrico as
coloca todas na fronteira); mesmo descontando essa faixa, a
concordância entre os métodos permanece.}

\textbf{O que fica para a leitura do documento.} A segunda fase não
altera nenhuma conclusão; ela redistribui os pesos. O \textbf{regime
operacional mais intenso e a eficiência produtiva maior das OSS são
traços robustos da rede}, visíveis em todos os métodos e em todas as
amostras, inclusive sem os conversores; o \textbf{TMP menor} replica
com exatidão em especificação hierárquica; o \textbf{custo} segue sem
diferença em qualquer método; e a \textbf{mortalidade menor}
concentra se inteiramente na trajetória dos cinco conversores, o
que a aproxima das ressalvas causais das
Seções~\ref{sec:efeitogestao} e~\ref{sec:inferenciafinal} e a afasta
de uma leitura de superioridade da rede OSS nesse desfecho.
